\section{Continue from the loader into the container}
Keep your completed loader. The attachment helper in
\path{packages/033-png/construction/png-collector-construction/attach.py} replaces only the two original
decoder stubs, adds a capacity control and appends the unfinished collector. It leaves
your loader instructions intact and refuses an existing output. Run it with
\texttt{--source} pointing to your current program and \texttt{--output} naming a new
working file. A changed decoder tail must be reconciled by hand; the helper refuses
to guess which of your functions it may discard.

The temporary decoder constructs the collector's five-field descriptor on its own
stack.

Set \texttt{collector\_capacity} to 46 bytes for the worked
image, within the actual compressed buffer capacity. Stop before
\texttt{collector\_call} to inspect the four arguments. Stop before
\texttt{collector\_return} to inspect the collector status and descriptor while that
frame is still alive. The supplied program still prints the development message after
successful collection, because it has not inflated or displayed the image. Both mode
settings follow this same temporary path; there is no vector operation to compare yet.

The collector has seventeen editable regions. Its supplied 128-byte frame saves
\texttt{\$ra}, \texttt{\$s0} through \texttt{\$s7}, and \texttt{\$fp}. It reserves a
native doubleword at frame offset 80 for the current payload length, at 88 for the type
value and at 104 for the computed CRC. These locations are construction aids, not part
of the public collector interface. Each helper is internal to this collector. Until
filled, a helper leaves through \texttt{collect\_missing}, which restores the owning
collector frame and returns development status four to its caller. Do not call an
unfinished helper as a standalone function.

Use \texttt{\$s0} for the current chunk address, \texttt{\$s1} for its remaining file
bytes, \texttt{\$s2} for the next compressed destination byte, \texttt{\$s3} for its
remaining capacity and \texttt{\$s4} for the descriptor. The construction uses
\texttt{\$s5} as a bit set: 1 means IHDR has been accepted, 2 means PLTE has been
accepted, 4 means the IDAT sequence has begun, and 8 means that sequence has ended.
The dimensions stay in \texttt{\$s6}/\texttt{\$s7}; the joined length stays in
\texttt{\$fp}. Preserve this state across helper calls. The completed-interface checks
also accept other internal arrangements; these checkpoints follow the illustrated one.

\Needspace{6\baselineskip}
Prepare each named observation from the file containing your current edits. The
collector preparation helper changes path and numeric controls, not instruction bodies.
Its fixture helper creates each input at a fresh path. Keep the loader input capacity at 1 MiB; the new control selects a smaller compressed
window in its separate allocation.
